\documentclass[11pt,a4paper]{article}
\usepackage[a4paper,margin=2.5cm]{geometry}
\usepackage[english]{babel}
\usepackage{amsmath}
\usepackage{amssymb}
\usepackage{amsthm}
\usepackage{array}
\usepackage{booktabs}
\usepackage{longtable}
\usepackage{seqsplit}
\usepackage{hyperref}

\newtheorem{theorem}{Theorem}[section]
\newtheorem{lemma}[theorem]{Lemma}
\newtheorem{proposition}[theorem]{Proposition}
\newtheorem{conjecture}[theorem]{Conjecture}
\theoremstyle{definition}
\newtheorem{definition}[theorem]{Definition}
\newtheorem{remark}[theorem]{Remark}
\newcommand{\coefficients}[1]{\texttt{\seqsplit{#1}}}

\title{Exact Ehrhart Numerators and Real-Rootedness for Multilayer Crown Order Polytopes: Tensor Networks and an Exact Theorem for $n=3$}
\author{Wendel León Alvarado}
\date{November 11, 2026}

\begin{document}

\maketitle

\begin{abstract}
We study Ehrhart $h^*$-polynomials of order polytopes associated with the multilayer crown posets $C_{k,n}$. A proposed reduction to a one-dimensional transfer trace is valid in height two but fails already for $C_{3,3}$: at dilation one, the exact lattice-point count is $28$, whereas the proposed trace gives $18$. We express the count as a cylindrical two-dimensional tensor-network contraction and reconstruct the Ehrhart numerator exactly over $\mathbb{Z}$. For fixed dilation, the implemented row transfer gives an integer recurrence in the layer index by Cayley--Hamilton; its state-space bound depends on the dilation. Exact CRT contractions and rational Sturm or reduced $\Gamma$-Sturm certificates establish simple, strictly negative real roots in 19 cases, through $C_{8,5}$ of degree 32. Exact symbolic reconstruction finds finite-range recurrences for $n=4$ and, for $n=5$, rules out two bounded-degree order-four models; specialization at $t=1$ has an exact scalar order-three recurrence with coefficients $(384,-10240,-32768)$. An exact six-state transfer proves the order-two recurrence for every $n=3$ Ehrhart numerator. A phase argument then proves that all $2k$ roots are simple and lie in $(-7-4\sqrt3,-7+4\sqrt3)$ for every $k\ge2$. The remaining real-rootedness and uniform recurrence questions concern circumferences $n\ge4$.
\end{abstract}

\small\noindent\textbf{Keywords:} Crown posets, Ehrhart polynomials, tensor networks, Sturm method, real-rootedness.
\normalsize

\section{Introduction \& Background}

For a finite poset $P$, its order polytope is
\begin{equation}
\mathcal{O}(P)=\{x\in[0,1]^P : x_u\le x_v\text{ whenever }u\le_P v\}.
\end{equation}
It is a full-dimensional lattice polytope of dimension $|P|$; its vertices are the characteristic vectors of upper sets of $P$ \cite{Stanley1986}.

Fix integers $k\ge2$ and $n\ge3$. The poset $C_{k,n}$ has elements $x_{j,i}$ for $1\le j\le k$ and $i\in\mathbb{Z}/n\mathbb{Z}$. Its cover relations are
\begin{equation}
x_{j,i}\prec x_{j+1,i},\qquad x_{j,i}\prec x_{j+1,i+1}
\quad (1\le j<k),
\end{equation}
where the second index is read cyclically. The rank of $x_{j,i}$ is $j-1$, and $C_{k,n}$ has $d=kn$ elements.

Let
\begin{equation}
L_{k,n}(m)=|m\mathcal{O}(C_{k,n})\cap\mathbb{Z}^{d}|,
\qquad
\sum_{m\ge0}L_{k,n}(m)t^m=\frac{H_{k,n}(t)}{(1-t)^{d+1}}.
\end{equation}
Thus $H_{k,n}(t)=h^*(\mathcal{O}(C_{k,n}),t)$ and $H_{k,n}(0)=1$.

Lundström and Saud Maia Leite formulated the real-rootedness conjecture
for the simple one-dimensional crown family in work originating in 2025
\cite{LundstromLeite2026}; Cao and Fu subsequently proved the height-two
case $k=2$ \cite{CaoFu2026}. The present computational study addresses the
multilayer families $C_{k,n}$ with $k\ge3$. Its contributions are the
failure of the proposed one-dimensional transfer trace beyond height two,
a cylindrical two-dimensional tensor-network contraction, exact
reconstruction of the Ehrhart numerators, and finite-data recurrence
discovery. The recurrence and root-location statements for arbitrary
parameters remain conjectural.

\subsection{Methodological correction: the one-dimensional trace}

For $N\ge1$, define the matrix $K_N=(\min(a,b))_{1\le a,b\le N}$. In height two, setting $N=m+1$ gives the transfer identity
\begin{equation}\label{eq:height-two-trace}
L_{2,n}(m)=\operatorname{tr}(K_N^n).
\end{equation}
Indeed, fix the values on the bottom layer. For adjacent values $a_{i-1},a_i$, the top-layer value at their common upper neighbor has $N-\max(a_{i-1},a_i)$ choices. Replacing each bottom value $a_i$ by $N-a_i$ turns this factor into the matrix entry $\min(N-a_{i-1},N-a_i)$. Summing around the cycle yields \eqref{eq:height-two-trace}.

\subsection{Failure of the proposed higher-layer extension}

The height-two derivation sums out one layer and leaves pairwise interactions between neighboring values on a single cycle. With three or more layers, the intermediate variables remain coupled through two-dimensional local constraints; replacing those constraints by the proposed single transfer matrix does not preserve the count. The following exact example disproves that proposed extension.

\begin{proposition}\label{prop:transfer-counterexample}
For $C_{3,3}$ at dilation $m=1$, the lattice-point count is
\[
L_{3,3}(1)=28,
\]
whereas the proposed one-dimensional formula gives $\operatorname{tr}(K_2^3)=18$.
\end{proposition}

\begin{proof}
At $m=1$, represent each layer by the subset of its three cyclic positions assigned value $1$. A subset $A$ in one layer forces the next layer to contain $A\cup(A+1)$. Counting the compatible layer choices gives $18+6+3+1=28$. In contrast,
\[
K_2=\begin{pmatrix}1&1\\1&2\end{pmatrix},
\qquad \operatorname{tr}(K_2^3)=18.
\]
The two counts differ, so the proposed one-dimensional transfer-trace identity is not valid uniformly for $k\ge3$.
\end{proof}

\begin{remark}
This counterexample rules out the stated transfer-trace reduction; it does not rule out other one-dimensional representations for special parameters, nor does it disprove Conjecture~\ref{conj:realrooted}.
\end{remark}

\section{2D Boundary MPS Tensor Network Engine}

\subsection{Local factor and exact lattice-point count}

For a fixed dilation $m\ge0$, each coordinate takes values in $\{0,\ldots,m\}$. Define the local tensor
\begin{equation}\label{eq:local-tensor}
T_{a,b,c}=
\begin{cases}
1,&c\ge\max(a,b),\\
0,&\text{otherwise},
\end{cases}
\qquad 0\le a,b,c\le m.
\end{equation}
Its three indices have link dimension $D=m+1$. The indices $a,b$ encode adjacent values in one layer, and $c$ encodes the value above them. Thus the local factor enforces the required order constraint $x_{j+1,i}\ge\max(x_{j,i},x_{j,i+1})$.

The complete cylindrical network is the factor-graph contraction
\begin{equation}\label{eq:tensor-count}
L_{k,n}(m)=
\sum_{\substack{x_{j,i}\in\{0,\ldots,m\}\\
1\le j\le k,\ i\in\mathbb{Z}/n\mathbb{Z}}}
\prod_{j=1}^{k-1}\prod_{i\in\mathbb{Z}/n\mathbb{Z}}
T_{x_{j,i},\,x_{j,i+1},\,x_{j+1,i}}.
\end{equation}
The cyclic index creates the periodic boundary in the circumferential direction. Each coordinate is summed once, and every factor in the product enforces one local cover constraint.

\subsection{Boundary contraction and numerical modes}

The tensor network is contracted by organizing its multilinear factors
along the discrete cylindrical boundary and selecting an efficient
hierarchy of index contractions. This is an exact contraction of the
entire network: no boundary-rank compression or singular-value
truncation is applied. Accordingly, ``Boundary MPS'' refers here to the
boundary organization of the tensor contraction, not to an approximate
matrix-product-state method.

The exact contraction is carried out over arbitrary-precision integers,
so \eqref{eq:tensor-count} involves no floating-point rounding. For
spectral exploration, a separate numerical evaluation uses real or
complex double-precision tensor arithmetic on the selected computing
device. Numerical values are not used as integer certificates.

\subsection{Layer transfer and its recurrence bound}

For fixed $m$, let $X_{m,n}=\{0,\ldots,m\}^n$ be the set of row
configurations and define the integer matrix
\[
B_{m,n}(x,y)=\prod_{i\in\mathbb Z/n\mathbb Z}
T_{x_i,x_{i+1},y_i},
\qquad x,y\in X_{m,n}.
\]
Summing over the intermediate rows in \eqref{eq:tensor-count} gives
\[
L_{k,n}(m)=\mathbf 1^{\mathsf T}B_{m,n}^{\,k-1}\mathbf 1.
\]
The cyclic shift of coordinates commutes with $B_{m,n}$, so the transfer
preserves the subspace fixed by $C_n$. Its rank is the number of cyclic
orbits of $(m+1)$-colorings of an $n$-cycle,
\[
M_{m,n}=\frac1n\sum_{r=0}^{n-1}(m+1)^{\gcd(n,r)}.
\]
The orbit-sum basis makes the restricted transfer an integer matrix.
If its characteristic polynomial is
$z^{r}+\sum_{j=0}^{r-1}c_jz^j$, then
Cayley--Hamilton yields
\[
L_{k+r,n}(m)+\sum_{j=0}^{r-1}c_jL_{k+j,n}(m)=0,
\qquad r\le M_{m,n}.
\]
Indeed, multiply the characteristic-polynomial identity for
$B_{m,n}|_{V_{m,n}^{C_n}}$ by
$\mathbf 1^{\mathsf T}B_{m,n}^{k-1}$ on the left and $\mathbf 1$ on the
right. The coefficient list $(c_0,\ldots,c_{r-1})$ lies in
$\mathbb Z^{r}$.
This proves an integer recurrence in $k$ for every fixed dilation.
The bound and its coefficients depend on $m$, however, so this transfer
does not by itself prove a recurrence for $H_{k,n}(t)$ with coefficients
in $\mathbb Z[t]$ and a bound depending only on $n$. In particular, the
implemented row transfer is not invariant under reflection: for
$n=3,m=1$, the transition $(1,0,0)\to(1,0,1)$ is allowed, but reflecting
both rows gives the forbidden transition $(1,0,0)\to(1,1,0)$. Thus the
cyclic orbit count of binary descent masks,
\[
N_{C_n}(n)=\frac1n\sum_{d\mid n}\varphi(d)2^{n/d},
\]
is not the state-space dimension of this implemented operator. A
separate P-partition transfer on descent masks would need an independent
matrix-element identity with $H_{k,n}(t)$ before its
Cayley--Hamilton bound could be applied. A uniform polynomial recurrence
for the Ehrhart numerators therefore remains a separate question.

\section{Exact Algebraic Reconstruction \& Verification Results}

\subsection{Exact reconstruction of the Ehrhart numerator}

Put $d=kn$. Ehrhart theory gives
\begin{equation}\label{eq:binomial-basis}
L_{k,n}(m)=\sum_{j=0}^{d}h_j\binom{m+d-j}{d},
\qquad H_{k,n}(t)=\sum_{j=0}^{d}h_jt^j.
\end{equation}
We evaluate \eqref{eq:tensor-count} for $m=0,1,\ldots,d$ and solve the
resulting triangular linear system by exact integer arithmetic. The
coefficients above the structural degree $d-k$ vanish. No
floating-point interpolation is used.

\subsection{Exact results}

Table~\ref{tab:exact} reports all 19 exact numerators obtained by CRT
contraction and binomial-basis reconstruction. In every case the
certificate counts as many negative roots as the degree and verifies
simplicity, so all roots are real, strictly negative, and simple. For
the even-degree palindromic cases $C_{5,5}$ through $C_{8,5}$, Sturm's
method is applied exactly to the reduced $\Gamma$ polynomial. The
coefficients are listed in descending order.

\small
\setlength{\tabcolsep}{2pt}
\begin{longtable}{@{}c c c c c >{\raggedright\arraybackslash}p{0.52\textwidth}@{}}
\caption{Exact CRT dataset and root certificates; coefficient lists are descending.}\label{tab:exact}\\
\toprule
$k$ & $n$ & Degree & Negative roots & Method & Coefficients of $H_{k,n}(t)$\\
\midrule
\endfirsthead
\caption[]{Exact CRT dataset (continued).}\\
\toprule
$k$ & $n$ & Degree & Negative roots & Method & Coefficients, descending\\
\midrule
\endhead
\midrule
\multicolumn{6}{r@{}}{\emph{Continued on next page}}\\
\endfoot
\bottomrule
\endlastfoot
2&3&4&4/4&CRT+Sturm&\coefficients{1,11,24,11,1}\\
3&3&6&6/6&CRT+Sturm&\coefficients{1,18,93,160,93,18,1}\\
4&3&8&8/8&CRT+Sturm&\coefficients{1,25,211,739,1120,739,211,25,1}\\
5&3&10&10/10&CRT+Sturm&\coefficients{1,32,378,2088,5757,8064,5757,2088,378,32,1}\\
6&3&12&12/12&CRT+Sturm&\coefficients{1,39,594,4550,19005,44547,59136,44547,19005,4550,594,39,1}\\
7&3&14&14/14&CRT+Sturm&\coefficients{1,46,859,8468,48263,165086,344061,439296,344061,165086,48263,8468,859,46,1}\\
8&3&16&16/16&CRT+Sturm&\coefficients{1,53,1173,14185,103331,473631,1393927,2657795,3294720,2657795,1393927,473631,103331,14185,1173,53,1}\\
3&4&9&9/9&CRT+Sturm&\coefficients{1,69,1204,7164,16906,16906,7164,1204,69,1}\\
4&4&12&12/12&CRT+Sturm&\coefficients{1,100,3058,37164,204639,553200,767068,553200,204639,37164,3058,100,1}\\
5&4&15&15/15&CRT+Sturm&\coefficients{1,131,5869,116355,1144241,6007079,17690249,30135467,30135467,17690249,6007079,1144241,116355,5869,131,1}\\
6&4&18&18/18&CRT+Sturm&\coefficients{1,162,9641,273808,4129580,35320040,179517436,560710448,1100896046,1376570220,1100896046,560710448,179517436,35320040,4129580,273808,9641,162,1}\\
7&4&21&21/21&CRT+Sturm&\coefficients{1,193,14374,539238,11319081,141272409,1094268784,5432987952,17714213382,38591690550,56806162004,56806162004,38591690550,17714213382,5432987952,1094268784,141272409,11319081,539238,14374,193,1}\\
2&5&8&8/8&CRT+Sturm&\coefficients{1,112,1983,9684,16120,9684,1983,112,1}\\
3&5&12&12/12&CRT+Sturm&\coefficients{1,228,11096,175640,1128655,3314308,4717744,3314308,1128655,175640,11096,228,1}\\
4&5&16&16/16&CRT+Sturm&\coefficients{1,349,32235,1090390,16336920,122193552,491902733,1115855805,1462319070,1115855805,491902733,122193552,16336920,1090390,32235,349,1}\\
5&5&20&20/20&CRT+$\Gamma$&\coefficients{1,470,67430,3915550,107636520,1560009346,12847278850,63317014410,193454540375,374892046080,466747051056,374892046080,193454540375,63317014410,12847278850,1560009346,107636520,3915550,67430,470,1}\\
6&5&24&24/24&CRT+$\Gamma$&\coefficients{1,591,117186,10199141,445096496,10693336966,151557685236,1333861396106,7570649136251,28485470286251,72457863450986,126239758904721,151779252595816,126239758904721,72457863450986,28485470286251,7570649136251,1333861396106,151557685236,10693336966,445096496,10199141,117186,591,1}\\
7&5&28&28/28&CRT+$\Gamma$&\coefficients{1,712,181578,21653944,1361039825,48903716880,1068736046140,14905774156760,137717986745175,867511817670400,3808653657683030,11848821044700960,26437684248135895,42655638313857400,50003623795120200,42655638313857400,26437684248135895,11848821044700960,3808653657683030,867511817670400,137717986745175,14905774156760,1068736046140,48903716880,1361039825,21653944,181578,712,1}\\
8&5&32&32/32&CRT+$\Gamma$&\coefficients{1,833,260611,40040390,3382607635,168934442784,5270888638227,107372156127054,1479487760034900,14185823530891700,96815201555014565,478811892910351720,1740289156615238005,4698922195864486730,9500786705237707205,14463987063668528165,16633278213336049990,14463987063668528165,9500786705237707205,4698922195864486730,1740289156615238005,478811892910351720,96815201555014565,14185823530891700,1479487760034900,107372156127054,5270888638227,168934442784,3382607635,40040390,260611,833,1}\\
\end{longtable}
\normalsize

\subsection{Sanity Checks \& Data Verification}

The 19 certified polynomials were checked against the exact coefficient
lists in the JSON artifacts. Each satisfies
\[
H_{k,n}(0)=1,\qquad
t^{k(n-1)}H_{k,n}(1/t)=H_{k,n}(t),\qquad
\deg H_{k,n}=k(n-1).
\]
The second identity is strict palindromicity (the Gorenstein reciprocity
of these Ehrhart numerators); the degree agrees with
Lemma~\ref{lem:structure}. Independently, the JSON Sturm certificates
record a negative-root count equal to the degree and a simplicity check
for each of the 19 cases. Thus these checks apply to the listed finite
dataset, not to the full infinite family.

\subsection{Rational Sturm certificates}

For a polynomial $H\in\mathbb{Z}[t]$, we form its Sturm sequence over $\mathbb{Q}$. Sign variations at $-\infty$ and at $0$ give the exact number of distinct roots in $(-\infty,0)$; this calculation involves rational polynomial arithmetic, not approximate root locations. Simplicity is checked independently by a gcd computation. For even-degree palindromic $H$, the optional reduced method uses $H(t)=(1+t)^d\Gamma(t/(1+t)^2)$ and Sturm's sequence of $\Gamma$. A full count of the degree in negative simple roots certifies the same conclusion for $H$.

Sturm sequences can also isolate individual roots: for rational endpoints $a<b$, the difference in sign variations at $a$ and $b$ counts the roots in $(a,b)$, and rational bisection can refine an interval until it contains exactly one root. The current JSON output records the exact polynomial and Sturm-sequence coefficients, endpoint sign lists, variation counts, root count, and simplicity/reciprocity checks; it does not report decimal approximations as certificates.

Exact rational Sturm certificates were computed for all 19 configurations:
$n=3$ for $k=2,\ldots,8$, $n=4$ for $k=3,\ldots,7$, and $n=5$ for
$k=2,\ldots,8$. The four cases $C_{5,5}$ through $C_{8,5}$
use the exact reduced-$\Gamma$ method. Each JSON certificate records the
degree check, endpoint signs, variation counts, root count, and
simplicity check. The 19 exact certificates are
structured in JSON format, named
\texttt{certificado\_k\_n.json}, and available in the project's
supplementary repository.
\footnote{Execution scripts, CLI usage, and raw JSON certificate files are available in the open-source repository associated with this paper.}

\subsection{Degree and reciprocity}\label{sec:structure}

\begin{lemma}\label{lem:structure}
For $k\ge2$ and $n\ge3$, $H_{k,n}(t)$ is monic of degree $kn-k$ and satisfies
\begin{equation}\label{eq:reciprocity}
H_{k,n}(t)=t^{kn-k}H_{k,n}(1/t).
\end{equation}
\end{lemma}

\begin{proof}
Let $P=C_{k,n}$ and $d=kn$. Lattice points of $m\mathcal{O}(P)$ are weakly order-preserving maps $P\to\{0,1,\ldots,m\}$. Interior lattice points are maps $f:P\to\{1,\ldots,m-1\}$ that are strictly increasing on every cover relation.

Every maximal chain has $k$ elements. An interior lattice point therefore requires at least $k$ available integer values, so the smallest dilation with an interior lattice point is $m=k+1$. At this dilation, each vertical chain $x_{1,i}\prec\cdots\prec x_{k,i}$ forces $f(x_{j,i})=j$; these values satisfy all cover relations, giving exactly one interior lattice point. The codegree is consequently $k+1$. The standard degree-codegree relation for Ehrhart numerators gives $\deg H_{k,n}=d+1-(k+1)=d-k$, and the top coefficient equals the number of interior lattice points at the codegree, namely one. Thus $H_{k,n}$ is monic.

For reciprocity, let $\rho(x_{j,i})=j-1$. If $f$ is an interior lattice point of $m\mathcal{O}(P)$, then
\[
g(x)=f(x)-\rho(x)-1
\]
is weakly order-preserving and takes values in $\{0,\ldots,m-k-1\}$. Conversely, every such weak map gives an interior lattice point by $f(x)=g(x)+\rho(x)+1$. Hence the interior Ehrhart count is $E^\circ(m)=E(m-k-1)$ for $m\ge k+1$, and it is zero for smaller positive $m$. If $F(t)=\sum_{m\ge0}E(m)t^m$, this says
\[
\sum_{m\ge1}E^\circ(m)t^m=t^{k+1}F(t).
\]
Ehrhart reciprocity gives $\sum_{m\ge1}E^\circ(m)t^m=(-1)^{d+1}F(1/t)$. Substituting $F(t)=H_{k,n}(t)/(1-t)^{d+1}$ and simplifying yields \eqref{eq:reciprocity}.
\end{proof}

\begin{remark}\label{rem:minusone}
If $s=kn-k$ is odd, equation~\eqref{eq:reciprocity} forces $H_{k,n}(-1)=0$. Under Conjecture~\ref{conj:realrooted}, this root is simple and the other roots occur in reciprocal pairs.
\end{remark}

\begin{proposition}\label{prop:k23}
For $k=2$ and $n=3$,
\[
H_{2,3}(t)=t^4+11t^3+24t^2+11t+1,
\]
and its four roots are distinct and negative.
\end{proposition}

\begin{proof}
Equation~\eqref{eq:height-two-trace} gives $L_{2,3}(m)=1,18,129,571,1884$ for $m=0,1,2,3,4$. Substitution into \eqref{eq:binomial-basis} gives the displayed polynomial. Dividing by $t^2$ and setting $z=t+t^{-1}$ reduces its root equation to $z^2+11z+22=0$. Its roots $(-11\pm\sqrt{33})/2$ are distinct and less than $-2$, so each gives two distinct negative reciprocal roots of $t+t^{-1}=z$.
\end{proof}

\section{Symbolic Recurrence Reconstruction}
\label{sec:recurrences}

Solving exact coefficient systems over $\mathbb{Q}[t]$ against the
tensor-network dataset gives the following finite-range propositions.

\begin{proposition}[Finite exact recurrence for $n=3$]
\label{prop:recurrence-n3}
For the certified levels $k=2,\ldots,8$, the exact recurrence
\begin{equation}\label{eq:recurrence-n3}
H_{k,3}(t)
=(t^2+6t+1)H_{k-1,3}(t)
+(t^3-2t^2+t)H_{k-2,3}(t)
\end{equation}
holds for every tested transition $k=4,\ldots,8$. The recurrence
operator was reconstructed uniquely in the specified polynomial search
space, and the coefficient system is overdetermined by the exact data.
This proposition concerns only those certified transitions.
\end{proposition}

\begin{proof}
The coefficients are determined over $\mathbb{Q}$ from the certified
polynomials, then substituted into each transition using exact integer
polynomial arithmetic. The identity holds at $k=4,5,6,7,8$. The
uniqueness claim is relative to the bounded-degree search used by the
analyzer; it is not a uniqueness assertion among unrestricted operators
of arbitrary degree.
\end{proof}

\begin{proposition}[Finite exact recurrence for $n=4$]
\label{prop:recurrence-n4}
For the certified levels $k=3,\ldots,7$, the exact recurrence
\begin{equation}\label{eq:recurrence-n4}
H_{k,4}(t)
=P_4(t)H_{k-1,4}(t)
+Q_4(t)H_{k-2,4}(t)
+R_4(t)H_{k-3,4}(t)
\end{equation}
holds for the tested transitions $k=6,7$, where
\begin{align*}
P_4(t)&=t^3+23t^2+23t+1,\\
Q_4(t)&=7t^5-28t^4-22t^3-28t^2+7t,\\
R_4(t)&=t^8-5t^7+9t^6-5t^5\\
&\qquad{}-5t^4+9t^3-5t^2+t.
\end{align*}
Within the specified degree bounds and reciprocal-polynomial ansatz,
the exact coefficient system is overdetermined and has a unique integral
solution. This proposition asserts the identity only at $k=6,7$.
\end{proposition}

\begin{proof}
The three operators were reconstructed over $\mathbb{Q}$ using reciprocal
spans $(3,6,9)$ and degree bounds $(4,7,10)$. Exact substitution into
$H_{6,4}$ and $H_{7,4}$ gives zero residual in both transitions. The
uniqueness statement is relative to this finite-dimensional ansatz.
\end{proof}

For $n=5$, an unrestricted order-four search was also tested without
symmetry or divisibility conditions. With degree bounds $(4,12,16,20)$,
the coefficient matrix has 56 columns and 99 equations; modulo
$p=2305843009213693951$, its rank and the augmented rank are $56$ and
$57$. With expanded bounds $(6,14,20,26)$, the corresponding values are
70 columns, 117 equations, and ranks $70$ and $71$. These modular ranks
certify inconsistency over $\mathbb Q$: the coefficient matrix has full
column rank, while the augmented matrix has strictly larger rank.
Accordingly, neither bounded-degree model admits an exact order-four
recurrence.

At $t=1$, the exact sequence has minimum scalar recurrence order three:
\begin{equation}\label{eq:scalar-n5}
H_{k,5}(1)=384H_{k-1,5}(1)
-10240H_{k-2,5}(1)-32768H_{k-3,5}(1),
\qquad k=5,6,7,8.
\end{equation}
The four certified transitions determine the three scalar coefficients
uniquely and exclude scalar orders one and two.

The recurrence operators were uniquely reconstructed over
$\mathbb{Q}[t]$ and certified on exact integer arithmetic. Here
``uniquely'' is relative to the bounded-degree search space for $n=3$
and, for $n=4$, to the stated degree bounds and reciprocal spans; it does
not assert uniqueness without these restrictions.

All operator coefficients lie in $\mathbb{Z}[t]$ and are reciprocal with
their indicated spans:
\[
\begin{aligned}
P_3(t)&=t^2P_3(1/t),&
Q_3(t)&=t^4Q_3(1/t),\\
P_4(t)&=t^3P_4(1/t),&
Q_4(t)&=t^6Q_4(1/t),&
R_4(t)&=t^9R_4(1/t),
\end{aligned}
\]
where $P_3(t)=t^2+6t+1$ and $Q_3(t)=t^3-2t^2+t$. The operators have order $n-1$ for the two circumferences investigated.
The $n=3$ identities were initially
checked for transitions $k=4,\ldots,8$; Proposition~\ref{prop:n3-transfer}
below proves the recurrence for every $k$. For $n=4$, the two independent
transitions $k=6,7$ determine a unique integral solution under reciprocal
spans $(3,6,9)$ and degree bounds $(4,7,10)$.

\begin{remark}
The coefficients shown here are obtained by exact reconstruction and
substitution into the certified polynomials. In particular, the
$t^3$ coefficient of $Q_3$ and the $t^8$ coefficient of $R_4$ are both
$+1$; changing either to $-1$ does not reproduce the dataset. These
reconstruction checks are finite; the transfer proof below establishes
the $n=3$ identity for every $k$, while the $n=4$ identity remains
verified only on the stated range.
\end{remark}

\subsection[BKW spectral analysis for n=4]{BKW spectral analysis for the order-three $n=4$ recurrence}
\label{sec:n4-bkw}

In this subsection, assume that the $n=4$ recurrence
\eqref{eq:recurrence-n4} holds for every $k\ge6$, with the exact initial
polynomials $H_{3,4},H_{4,4},H_{5,4}$ from the certificates. The
recurrence itself has so far been verified only at $k=6,7$.
Let
\[
\chi_4(\lambda;t)=\lambda^3-P_4(t)\lambda^2-Q_4(t)\lambda-R_4(t),
\qquad R_4(t)=t(t-1)^6(t+1).
\]
Its discriminant is
\[
\Delta_4(t)=-4t\,q_4(t)\,q_6(t)^2,
\]
where
\[
q_4(t)=t^4+76t^3-282t^2+76t+1,\qquad
q_6(t)=t^6+7t^5+27t^4+58t^3+27t^2+7t+1.
\]
Writing $z=t+t^{-1}$ gives $q_4(t)/t^2=z^2+76z-284$. Its two
negative roots are
\[
\beta_\pm=\frac{z_0\pm\sqrt{z_0^2-4}}2,\qquad
z_0=-38-24\sqrt3,\qquad \beta_+\beta_-=1,
\]
with $\beta_-\simeq-79.5566497221$ and
$\beta_+\simeq-0.01256965953$. Set
$I_4=(\beta_-,\beta_+)$.

\begin{theorem}[Negative-real BKW accumulation set for $n=4$]
\label{thm:n4-bkw}
For the sequence generated by the order-three recurrence from
$H_{3,4},H_{4,4},H_{5,4}$, the set of accumulation points of its zeros
on the negative real axis is exactly
\[
[\beta_-,\beta_+].
\]
For $t\in I_4$ away from the zeros of $q_6$, the two nonreal conjugate
characteristic roots have equal modulus and strictly dominate the third
root. At the zeros of $q_6$, this dominant pair coalesces into a
repeated real root whose modulus still strictly exceeds that of the
third root. For
$t\in(-\infty,\beta_-)\cup(\beta_+,0)$, the three characteristic roots
are real and have pairwise distinct absolute values.
\end{theorem}

\begin{proof}
For $t<0$, the sign of $\Delta_4(t)$ is the sign of $q_4(t)$ away from
the zeros of $q_6$. Thus, in $I_4$ there is one real root and a
conjugate pair except at the zeros of $q_6$. To locate those exceptional
points, put $z=t+t^{-1}$; then
$q_6(t)/t^3=z^3+7z^2+24z+44$. This cubic is strictly increasing on
$\mathbb R$, is negative at $z_0$ and positive at $-2$, so its unique
root in $(z_0,-2)$ yields two reciprocal negative roots of $q_6$ in
$I_4$. At each, the characteristic cubic has a repeated dominant real
root. Outside the closure of $I_4$ there are three distinct real roots.
A tie between the modulus of the conjugate pair and that of the
real root would force
$P_4(t)^3R_4(t)+Q_4(t)^3=0$. Exact Sturm calculation shows that this
polynomial has no negative zero. At $t=-1$,
\[
\chi_4(\lambda;-1)=\lambda(\lambda^2+48),
\]
so the conjugate pair is strictly dominant there as well. Continuity
then gives the asserted strict dominance throughout $I_4$, including
the isolated repeated-eigenvalue points where $q_6(t)=0$.

Outside $I_4$, equality in absolute value between two distinct real
roots would require them to be opposites, which implies
$P_4Q_4+R_4=0$. Its exact factorization is
\[
P_4Q_4+R_4
=8t(t+1)(t^2-5t+1)(t^4+20t^3+22t^2+20t+1).
\]
All its negative roots lie in $I_4$; exact root isolation at
$t=-100$ and $t=-1/1000$ fixes the strict ordering on the two exterior
components.

For clarity, the proposed condition
$P_4^3R_4-Q_4^3\ne0$ throughout $I_4$ is false: that polynomial has
one root in each of $(-50,-49)$ and $(-21/1000,-1/50)$, both contained
in $I_4$. This condition is not needed for the dominance argument;
the relevant no-crossing condition is
$P_4^3R_4+Q_4^3\ne0$.

The numerator of the generating function for the recurrence is
\[
N(\lambda,t)=\lambda^2H_{3,4}
+\lambda(H_{4,4}-P_4H_{3,4})
+H_{5,4}-P_4H_{4,4}-Q_4H_{3,4}.
\]
The exact resultant is
\[
\operatorname{Res}_\lambda(\chi_4,N)
=16t^5(t-1)^{20}(t+1)^3q_6(t)^2.
\]
Consequently, away from the isolated repeated-eigenvalue points, the
dominant-mode coefficients are nonzero on the interior of $I_4$ and
the unique dominant coefficient is nonzero on the exterior. At
$t=-1$, $N(\lambda,-1)=1536\lambda$, so neither dominant coefficient
vanishes. the Beraha--Kahane--Weiss theorem \cite{BerahaKahaneWeiss1975} therefore gives the stated
accumulation set on $(-\infty,0)$; the closure includes the two
endpoints.
\end{proof}

\begin{remark}
The BKW accumulation theorem does not imply that every zero is real or
simple for all sufficiently large $k$. Exact Sturm checks show that
the recurrence-generated polynomials have $3k$ simple roots in
$(-79,-1/79)\subset I_4$ for $k=3,\ldots,12$, but this finite range is
not a proof for all large $k$. To make the assertion unconditional for
the actual $H_{k,4}$, one must first prove the order-three recurrence
for every $k$, and then establish a uniform real-zero count (for
example, a sign-change argument producing $3k$ distinct zeros).
No analogous recurrence or spectral analysis has been established here
for $n\ge5$.
\end{remark}

\subsection[Stage 1: exact transfer and real-rootedness for n=3]{Stage 1: an exact transfer and real-rootedness for $n=3$}
\label{sec:n3-real-rooted}

\begin{proposition}[Exact six-state transfer for $n=3$]
\label{prop:n3-transfer}
For the naturally labelled cylindrical poset, with
$L(j,i)=3j+i$ for $j\ge0$ and $i\in\{0,1,2\}$, its Ehrhart numerator
$H_{k,3}$ satisfies, for every $k\ge2$,
\[
H_{k,3}(t)=P_3(t)H_{k-1,3}(t)+Q_3(t)H_{k-2,3}(t),
\qquad P_3(t)=t^2+6t+1,\quad Q_3(t)=t(t-1)^2,
\]
with $H_{0,3}(t)=1$ and $H_{1,3}(t)=t^2+4t+1$.
\end{proposition}

\begin{proof}
For a naturally labelled poset, the Ehrhart $h^*$-polynomial is the
descent enumerator of its linear extensions \cite{Stanley1986}. At a
prefix of $3r$ elements, an order ideal is either $r$ complete layers or
$r-1$ complete layers followed by two partial layers of sizes two and one.
Indeed, once a layer is not full, no later layer can be full; a partial
layer of size two can be followed by at most one element, while one
element cannot be followed by a nonempty layer. The only possible
two-layer tails, written as bitmasks of vertices, are $(3,2)$, $(5,1)$,
and $(6,4)$. Remembering the last element gives the following six states:
\[
\begin{array}{c|c|c}
s&\text{partial-layer masks}&\text{last vertex}\\ \hline
1&(3,2)&(r,1)\\
2&(5,1)&(r,0)\\
3&(6,4)&(r,2)\\
4&\text{$r$ full layers}&(r-1,0)\\
5&\text{$r$ full layers}&(r-1,1)\\
6&\text{$r$ full layers}&(r-1,2)
\end{array}
\]
Here mask bit $i$ records whether vertex $i$ is present. Enumerating
the legal ways to append the next three elements, with a factor $t$
for each new descent, gives the column-vector transfer
\[
M(t)=
\begin{pmatrix}
t&t&0&t+1&t+1&t+1\\
0&t&t&t+1&t+1&t+1\\
t&0&t&t+1&t+1&t+1\\
t^2&0&t^2&t^2+t&t^2+t&t^2+t\\
0&t^2&t&2t&2t&2t\\
t&t&0&t+1&t+1&t+1
\end{pmatrix}.
\]
The exact initial prefix weights and terminal completion weights are
\[
\mathbf v=
\begin{pmatrix}
(t+1)P_3\\(t+1)P_3\\(t+1)P_3\\
t(t+1)P_3\\t(t+3)(3t+1)\\(t+1)P_3
\end{pmatrix},
\qquad
\alpha(t)=t(t+1),\quad \beta(t)=t^2+4t+1,\quad
\mathbf g=(\alpha,\alpha,\alpha,\beta,\beta,\beta).
\]
Thus $H_{k,3}(t)=\mathbf gM(t)^{k-3}\mathbf v$ for $k\ge3$.
The transfer matrix and its boundary vectors are derived by exact
enumeration in \texttt{TENSORES 2D/verify\_n3\_symbolic.py}.

Writing $\chi_3(\lambda;t)=\lambda^2-P_3(t)\lambda-Q_3(t)$, direct
symbolic calculation gives
\[
\det(\lambda I-M)=
\lambda^2(\lambda^2-\lambda t+t^2)\chi_3(\lambda;t),
\]
and the scalar resolvent is
\[
\mathbf g(\lambda I-M)^{-1}\mathbf v
=\frac{\lambda A(t)+B(t)}{\chi_3(\lambda;t)},
\]
where
\begin{align*}
A(t)&=t^6+18t^5+93t^4+160t^3+93t^2+18t+1,\\
B(t)&=t^7+9t^6+3t^5-26t^4+3t^3+9t^2+t.
\end{align*}
The numerator is coprime to $\chi_3$ over $\mathbb Q(t)$, so the scalar
transfer sequence has minimal denominator $\chi_3$. Its coefficients
therefore satisfy the asserted recurrence for $k\ge5$. Direct exact
enumeration gives $H_{0,3}=1$, $H_{1,3}=t^2+4t+1$ and verifies the
remaining initial cases $k=2,3,4$, proving the recurrence for every
$k\ge2$.

The six-state transfer is not $C_3$-equivariant: rotating the vertices
does not preserve descent weights for the natural labeling. Accordingly,
the proof uses the full six-state transfer and its scalar resolvent, not
a quotient by cyclic symmetry.
\end{proof}

\begin{theorem}[Real-rootedness and negative simple roots for $n=3$]
\label{thm:n3-real-rooted}
For every $k\ge2$, the actual Ehrhart numerator $H_{k,3}$ satisfies
\[
H_{k,3}(t)=P_3(t)H_{k-1,3}(t)+Q_3(t)H_{k-2,3}(t),
\quad
P_3(t)=t^2+6t+1,\quad Q_3(t)=t(t-1)^2,
\]
with $H_{0,3}(t)=1$ and $H_{1,3}(t)=t^2+4t+1$. It is monic of degree
$2k$ and has exactly $2k$ distinct, strictly negative real roots, all
lying in
\[
I_3=(-7-4\sqrt3,\,-7+4\sqrt3)\subset(-\infty,0).
\]
\end{theorem}

\begin{proof}
Proposition~\ref{prop:n3-transfer} establishes the recurrence and initial
values for the actual Ehrhart numerators. The characteristic
discriminant factors as
\[
\Delta_3(t)=P_3(t)^2+4Q_3(t)
=(t+1)^2(t^2+14t+1).
\]
Writing $A(t)=-(t^2+14t+1)$, we have $A(t)>0$ on $I_3$ and
\[
r(t)=(1-t)\sqrt{-t}>0.
\]
The characteristic roots are $r(t)e^{\pm i\theta(t)}$, where the
continuous choice $\theta(t)\in(0,2\pi)$ on $I_3$ is specified by
\[
\cos\theta(t)=\frac{t^2+6t+1}{2(1-t)\sqrt{-t}},\qquad
\sin\theta(t)=\frac{(t+1)\sqrt{A(t)}}{2(1-t)\sqrt{-t}}.
\]
Since $2H_{1,3}(t)-P_3(t)=(t+1)^2$, solving for the Binet coefficients
gives, away from the repeated-root point $t=-1$,
\[
c_\pm(t)=\frac12\mp\frac{i(t+1)}{2\sqrt{A(t)}}.
\]
The resulting expression extends continuously through $t=-1$, where
the characteristic root is repeated. The magnitude and argument of
$c_+$ are
\[
2|c_+(t)|=\sqrt{\frac{-12t}{A(t)}},\qquad
\phi(t)=\arg c_+(t)
=\arctan\!\left(\frac{-(t+1)}{\sqrt{A(t)}}\right)
\in\left(-\frac\pi2,\frac\pi2\right).
\]
Consequently, throughout $I_3$,
\begin{equation}\label{eq:n3-binet-phase}
H_{k,3}(t)=r(t)^k
\sqrt{\frac{-12t}{A(t)}}\,
\cos\bigl(\Psi_k(t)\bigr),
\qquad \Psi_k(t)=k\theta(t)+\phi(t).
\end{equation}
At $t=-1$, this identity holds by continuity (and gives
$H_{k,3}(-1)=(-2)^k$).

Differentiating the phase yields
\[
\Psi_k'(t)=
\frac{-(k+1)t^2+(10k+2)t-(k+1)}
{2t(t-1)\sqrt{A(t)}}.
\]
For $t\in I_3$, the denominator is positive, while the numerator is
$-(k+1)(t^2+1)+(10k+2)t<0$ for every $k\ge1$. Hence $\Psi_k$ is
strictly decreasing. At the two endpoints,
\[
\lim_{t\to(-7-4\sqrt3)^+}\Psi_k(t)
=2k\pi+\frac\pi2,\qquad
\lim_{t\to(-7+4\sqrt3)^-}\Psi_k(t)=-\frac\pi2.
\]
Therefore, for each $m=0,1,\ldots,2k-1$, there is exactly one
$t\in I_3$ with $\Psi_k(t)=(m+\tfrac12)\pi$. Equation
\eqref{eq:n3-binet-phase} makes each such $t$ a root. The prefactor is
strictly positive and $\Psi_k'(t)\ne0$, so differentiation of
\eqref{eq:n3-binet-phase} shows each root is simple. These are $2k$
distinct roots. Finally, induction in the recurrence shows that
$H_{k,3}$ is monic of degree $2k$: the term $P_3H_{k-1,3}$ has degree
$2k$ and leading coefficient one, while $Q_3H_{k-2,3}$ has degree
$2k-1$. Thus the displayed roots account for all roots.
\end{proof}

\section{Real-Rootedness \& Interlacing Implications}
\label{sec:root-implications}

\begin{conjecture}[Higher-layer negative real-rootedness]
\label{conj:realrooted}
For every $k\ge3$ and $n\ge4$, all zeros of $H_{k,n}(t)$ are negative
real numbers and are simple.
\end{conjecture}

\begin{conjecture}[Cylindrical recurrence of order $n-1$]
\label{conj:universal-recurrence}
For each fixed circumference $n\ge4$, there are polynomials
$A_{1,n}(t),\ldots,A_{n-1,n}(t)\in\mathbb{Z}[t]$, independent of $k$,
such that the recurrence holds for every $k\ge n$, where all lagged
polynomials $H_{k-j,n}$ are defined:
\[
H_{k,n}(t)=\sum_{j=1}^{n-1}A_{j,n}(t)H_{k-j,n}(t).
\]
\end{conjecture}

The six-state transfer proves the order-two recurrence for every
$H_{k,3}$. For $n=4$, the order-three identity is verified only on the
stated finite range. For $n=5$, the order-four searches above exclude
the two specified degree-bounded models, not every possible polynomial
operator of arbitrary degree. The fixed-dilation transfer theorem proves
recurrences for $L_{k,n}(m)$, but does not establish
Conjecture~\ref{conj:universal-recurrence} for $n\ge4$.
Theorem~\ref{thm:n3-real-rooted} settles real-rootedness unconditionally
for $n=3$; the remaining real-rootedness conjecture concerns $n\ge4$.

\subsection[Characteristic discriminant for the exact n=3 recurrence]{Characteristic discriminant for the exact $n=3$ recurrence}

With the convention $H_{k,3}=P_3H_{k-1,3}+Q_3H_{k-2,3}$, the
characteristic equation in the layer index is
$\lambda^2-P_3(t)\lambda-Q_3(t)=0$. Its discriminant is therefore
\[
\Delta_3(t)=P_3(t)^2+4Q_3(t)
=t^4+16t^3+30t^2+16t+1.
\]
For $t<0$, write $x=-t>0$ and $z=x+x^{-1}\ge2$. Then
\[
\frac{\Delta_3(-x)}{x^2}=z^2-16z+28=(z-2)(z-14).
\]
Consequently, $\Delta_3(t)$ vanishes at $z=2$ and $z=14$, is positive
when $z>14$, and is negative when $2<z<14$. Thus the characteristic
multipliers are not real for every negative $t$; this discriminant does
not by itself imply stability or real-rootedness of the polynomial
sequence. The Sturm certificates establish the latter only for the
individual cases in Table~\ref{tab:exact}.

The alternative expression
$P_3(t)^2-4(-t^3-2t^2+t)=t^4+16t^3+46t^2+8t+1$
is not positive for every $t<0$: at $t=-4$ it equals $-63$. Moreover,
the proposed $Q_3(t)=-t^3-2t^2+t$ differs from the verified operator
$t^3-2t^2+t$ by $-2t^3$, so it does not reproduce the certified
recurrences. Similarly, changing the leading coefficient of $R_4$ from
$+1$ to $-1$ changes the recurrence residual by a nonzero multiple of
$2t^8H_{k-3,4}$. The coefficients in Section~\ref{sec:recurrences} are
retained because they reproduce the exact data.

\subsection{Interlacing status}

The Sturm certificates prove that each listed polynomial has only simple
negative real roots. They do not establish strict interlacing between
different values of $k$. The degrees increase by $n-1$ in each fixed-$n$
family, so any interlacing claim requires a separately specified
generalized interlacing criterion and an independent proof; no such
general theorem is inferred from the current finite dataset.

\section{Concluding Remarks \& Future Work}

The height-two transfer trace is justified by summing out a single layer, but the explicit $C_{3,3}$ count shows that the proposed one-dimensional extension does not preserve lattice-point counts for higher layers. The cylindrical tensor network retains the full local coupling, and its fixed-dilation row transfer proves integer recurrences for the lattice-point counts. For $n=3$, the exact six-state transfer proves an order-two recurrence for every Ehrhart numerator; the phase argument then gives $2k$ simple roots in $I_3$ for all $k\ge2$. Exact CRT contraction and binomial-basis reconstruction yield reproducible integer polynomials; Sturm certificates verify simple negative real-rootedness for all 19 finite cases in Table~\ref{tab:exact}, through $C_{8,5}$. Thus the remaining real-rootedness conjecture concerns $n\ge4$. For $n=5$, two bounded-degree order-four systems are inconsistent, while the exact scalar specialization is \eqref{eq:scalar-n5}. A uniform polynomial recurrence for $n\ge4$ and higher-layer real-rootedness remain open. Future work includes constructing fixed-dimensional transfers for higher circumferences and proving root-location or generalized interlacing results uniformly in $k,n$.

The full-layer step raises the degree by $n-1$, while the finite tests of the transfer at $t=-s$ found neither total nonnegativity nor a diagonal sign conjugation to a nonnegative matrix. A vertex-by-vertex interpolation through order ideals, whose height profiles satisfy $h_i\le h_{i-1\pmod n}+1$, instead gives
$C_{k,n}=P_0\subset P_1\subset\cdots\subset P_n=C_{k+1,n}$ and the exact decomposition
$H_I(t)=\sum_{m\in\max(I)}H_I^{(m)}(t)$ by last maximal element. The maximal-deletion identity
\[
H_I^{(m)}(t)=\sum_{w\in\max(I\setminus\{m\})}
t^{\mathbf{1}[w>m]}H_{I\setminus\{m\}}^{(w)}(t)
\]
was verified on all tested ideals. The directed interlacing sequence also held at every step of the five tested intermediate chains for $n\in\{3,4,5\}$; this finite evidence suggests a natural induction program for the conjecture at all $n\ge4$, but does not establish such an induction. In particular, the exhaustive ideal test for $C_{3,5}$ had two exceptions to complete directed interlacing.

\subsection{Data \& Code Availability}

The source code is in the sibling project directory \texttt{TENSORES 2D/}:
the tensor-network engine and exact reconstruction modules are in
\texttt{core/}, with command-line entry point \texttt{main.py}. The 19
JSON certificates accompanying Table~\ref{tab:exact} are distributed
alongside this article. Sturm verification is carried out by
\texttt{core/root\_verifier.py}; recurrence and interlacing computations
are carried out by \texttt{core/interlacing\_analyzer.py}. The
certificates record the polynomial coefficients and exact Sturm data
needed to reproduce the finite checks reported here. The exact $n=3$
transfer proof is independently checked by
\texttt{verify\_n3\_symbolic.py}; the $n=4$ spectral calculations are
checked by \texttt{verify\_n4\_bkw.py}.

\begin{thebibliography}{9}
\bibitem{LundstromLeite2026}
T.~Lundstr\"om and L.~Saud Maia Leite,
\emph{Order polytopes of crown posets},
European Journal of Combinatorics \textbf{133} (2026), 104304.
\href{https://doi.org/10.1016/j.ejc.2025.104304}{doi:10.1016/j.ejc.2025.104304}.

\bibitem{Stanley1986}
R.~P.~Stanley,
\emph{Two poset polytopes},
Discrete \& Computational Geometry \textbf{1} (1986), 9--23.

\bibitem{BerahaKahaneWeiss1975}
{\raggedright
S.~Beraha, J.~Kahane, and N.~J.~Weiss,
\emph{Limits of zeros of recursively defined polynomials},
Proceedings of the National Academy of Sciences \textbf{72} (1975), 4209--4211.
\href{https://doi.org/10.1073/pnas.72.11.4209}{doi:10.1073/pnas.72.11.4209}.\par}

\bibitem{CaoFu2026}
Y.~Cao and H.~Fu,
\emph{Real-rooted Ehrhart $h^*$-polynomials of crown posets and wheel and whirl matroids},
preprint (2026), arXiv:2610.40496.
\href{https://arxiv.org/abs/2610.40496}{arXiv:2610.40496}.
\end{thebibliography}

\end{document}
